\begin{pastexam}[label={px:nsysu-111-osds-02}]{中山資工所甲組 111 學年度 第 2 題}
（問答）

(10\%) Suppose four symbols A, B, C, D are input into a stack sequentially. Four PUSH operations
and four POP operations are used for the stack. There are some sequences are impossible output
from the all possible combinations of PUSH and POP. Please give the number of each impossible
case with AXXX, BXXX, CXXX, DXXX separately.
\end{pastexam}

\begin{pastsoln}
答案：AXXX 有 $1$ 種不可能（ADBC），BXXX 有 $1$ 種（BDAC），CXXX 有 $3$ 種（CABD、CADB、CDAB），
DXXX 有 $5$ 種（DABC、DACB、DBAC、DBCA、DCAB），一共 $10$ 種。

\textbf{判斷的工具}。 把 A、B、C、D 看成 $1, 2, 3, 4$。 資結書定理~6.1：依序進堆疊時，輸出
順序做得到，若且唯若輸出裡找不到「先出現一個大的 $z$，之後依序出現 $x$、$y$，而
$x < y < z$」（俗稱 3-1-2 型）。 原因是 $z$ 出來的時候，比它小的 $x$、$y$ 都已經進了堆疊，
$x$ 在 $y$ 下面，只能 $y$ 先出來。 四個符號的排列共 $4! = 24$ 種，做得到的有卡特蘭數
$C_4 = \frac{1}{5}\binom{8}{4} = 14$ 種（資結書定理~6.2），所以不可能的有 $24 - 14 = 10$ 種。
下面依第一個輸出的符號分開數。

\textbf{AXXX}：A 先 push 再立刻 pop，堆疊變空，剩下 B、C、D 就是三個符號依序進堆疊的問題。
三個符號的 $6$ 種排列裡只有 $3$-$1$-$2$ 型做不到，也就是 D、B、C（D 出來時 B、C 都在堆疊裡，
C 在上面）。 不可能：\textbf{ADBC}，$1$ 種。

\textbf{BXXX}：push A、push B、pop B，此時堆疊裡有 A，還沒進的是 C、D。 後三個是 A、C、D 的
$6$ 種排列，只有 D 在最前面而後面是 A、C 的那一種違規：D 出來時 C 已經壓在 A 上面，A 不能比 C
早出來。 不可能：\textbf{BDAC}，$1$ 種。 其他五種都做得到，例如 BCDA（push C、pop C、push D、
pop D、pop A）、BDCA（push C、push D、pop、pop、pop）。

\textbf{CXXX}：要先輸出 C，A、B、C 都要 push，pop C 以後堆疊由底到頂是 A、B，還沒進的只有 D。
B 在 A 上面，所以後三個裡 \textbf{B 一定在 A 前面}，D 插在哪裡都可以（D 隨時可以 push 再 pop）。
A、B、D 的 $6$ 種排列裡 B 在 A 前面的有 BAD、BDA、DBA 三種，做得到；其餘三種 A 在 B 前面，
做不到。 不可能：\textbf{CABD、CADB、CDAB}，$3$ 種。

\textbf{DXXX}：先輸出 D，四個都要 push，堆疊由底到頂是 A、B、C，之後只能依序 pop 出 C、B、A。
唯一做得到的是 DCBA，其他 $3! - 1 = 5$ 種都做不到：\textbf{DABC、DACB、DBAC、DBCA、DCAB}。

\textbf{核對總數}：$1 + 1 + 3 + 5 = 10$，和 $24 - 14 = 10$ 相符。 每一種不可能的序列裡都能找到
3-1-2 型，例如 ADBC 的 D、B、C，CABD 的 C、A、B，DBCA 的 D、B、C。
\end{pastsoln}
